\subsection{Motivazione}

La soluzione numerica di una PDE mediante il metodo degli elementi finiti dipende non soltanto dal numero di gradi di libertà, ma anche dalla loro distribuzione spaziale. Quando la soluzione presenta fronti, picchi o regioni con variazioni elevate e localizzate, una mesh uniforme può allocare una quota non ottimale dei gradi di libertà: alcune regioni risultano sovra-risolte, mentre altre non descrivono adeguatamente le strutture della soluzione. L'adattività di mesh affronta questo problema modificando la discretizzazione in risposta alla soluzione numerica, a indicatori di errore o a informazioni fisiche disponibili \citep{huang2011adaptive,bangerth2003adaptive,cao2001an}.

Le principali famiglie di adattività si distinguono per il tipo di modifica apportata alla discretizzazione. Nell'$h$-refinement si modificano le dimensioni degli elementi mediante raffinamento o coarsening; nel $p$-refinement si modifica l'ordine polinomiale; nell'$r$-adaptation si mantengono invariati il numero di nodi e la connettività, modificando invece le coordinate dei nodi. Quest'ultima alternativa è interessante quando si desidera redistribuire una quantità prefissata di gradi di libertà senza introdurre nuovi elementi o modificare la struttura topologica della mesh \citep{huang2011adaptive,hyman1988moving}.

Nei problemi transitori, la rilocazione deve essere integrata con l'evoluzione dello stato discreto. Una modifica istantanea della geometria richiede infatti un operatore di trasferimento tra spazi agli elementi finiti differenti; una successione di trasferimenti può introdurre errori addizionali e aumentare il costo complessivo. Un approccio Arbitrary Lagrangian--Eulerian (\ALE) descrive invece il moto della griglia mediante una velocità di mesh e consente di incorporare la variazione geometrica direttamente nella formulazione temporale \citep{hirt1974an,donea2004arbitrary,farhat2001the}.

\subsection{Rilocazione di mesh e metodi basati sull'apprendimento}

I metodi classici di moving mesh determinano la distribuzione nodale attraverso equazioni ausiliarie, funzionali variazionali o monitor function costruite a partire dalla soluzione o da sue derivate \citep{huang2001practical,cao2001an,huang2011adaptive}. Queste strategie sono consolidate, ma la risoluzione del problema di mesh movement può diventare un costo significativo quando la griglia viene aggiornata frequentemente. Inoltre, la qualità di una configurazione geometricamente desiderabile non coincide necessariamente con la riduzione dell'errore della PDE discretizzata.

Le reti neurali su grafi forniscono una rappresentazione adatta a mesh non regolari: i nodi e gli elementi possono essere descritti mediante un grafo, mentre il message passing consente di costruire predizioni dipendenti dal contesto locale della discretizzazione. La letteratura ha già esplorato il movimento di mesh mediante reti neurali. M2N propone reti di movimento per solver PDE, includendo casi dipendenti dal tempo e strategie per ridurre il rischio di mesh tangling \citep{song2022m2n}. UM2N studia una rete di movimento più generale e la sua applicazione a differenti strutture di mesh e famiglie di problemi \citep{zhang2024towards}. G-Adaptivity propone una GNN addestrata minimizzando direttamente un errore della soluzione agli elementi finiti, e costituisce un riferimento particolarmente rilevante perché collega esplicitamente il movimento della mesh alla qualità della soluzione PDE \citep{rowbottom2025gadaptivity}. Un ulteriore studio recente affronta la $r$-adattività neurale per problemi ellittici parametrici autoaggiunti, utilizzando un funzionale di tipo Ritz \citep{aballay2026radaptive}.

Questi contributi escludono una formulazione del problema in termini di ``prima GNN per il movimento di mesh'' o di assenza generale di metodi neurali per l'adattività. Il presente lavoro non rivendica tale priorità. La questione affrontata è più circoscritta: valutare, su un problema transiente bidimensionale di convezione--diffusione, una pipeline che accoppia (i) una GNN che produce spostamenti nodali a partire dallo stato FEM corrente, (ii) un filtro deterministico che controlla la traiettoria geometrica proposta e (iii) un'integrazione ALE nativa basata su deal.II, verificando contemporaneamente accuratezza PDE, validità geometrica e costo computazionale.

\subsection{Problema aperto e obiettivo del lavoro}

La difficoltà metodologica non consiste soltanto nel predire una nuova configurazione nodale. In una simulazione transiente, la predizione deve essere compatibile con lo stato numerico corrente; il movimento deve preservare l'ammissibilità della mesh lungo il percorso, non soltanto al termine dello spostamento; e la discretizzazione mobile deve trattare coerentemente la variazione geometrica nel tempo. Infine, la qualità della strategia deve essere misurata sull'errore della PDE e non soltanto su un indicatore di interpolazione o su un obiettivo geometrico.

Il framework qui studiato affronta questi aspetti come una pipeline integrata. La GNN produce una proposta di spostamento; un filtro deterministico determina un fattore globale di applicazione compatibile con una soglia sul Jacobiano degli elementi triangolari; la traiettoria filtrata viene distribuita su una finestra temporale e utilizzata dalla formulazione ALE. Le operazioni di mesh, FEM, ALE, sicurezza geometrica e valutazione degli errori sono implementate nel nucleo C++ basato su deal.II, mentre Python gestisce l'inferenza neurale, i dataset e l'orchestrazione degli esperimenti \citep{arndt2021dealii}.

L'ipotesi sperimentale è che una GNN addestrata a rappresentare l'evoluzione temporale della mesh possa produrre configurazioni più utili per una PDE transiente rispetto a una GNN addestrata su spostamenti istantanei. Tale ipotesi viene verificata mediante un confronto appaiato con una mesh fissa e con il modello statico Stage A. Il risultato è valutato su casi held-out distinti in un regime interpolativo e uno estrapolativo rispetto agli intervalli dei parametri utilizzati in addestramento.

\subsection{Contributi e perimetro delle conclusioni}

I contributi del lavoro sono i seguenti:
\begin{enumerate}
    \item una pipeline di $r$-adattività bidimensionale basata su grafo, con topologia fissa e predizione dello spostamento nodale da feature derivate dallo stato FEM corrente;
    \item un filtro geometrico deterministico basato sull'analisi quadratica del Jacobiano lungo una traiettoria lineare di spostamento, con verifica numerica della sua applicazione;
    \item l'integrazione del movimento filtrato in un solutore ALE e la verifica della compatibilità geometrica discreta mediante test di free-stream e residui D-GCL;
    \item una valutazione progressiva dei componenti, comprendente target di rilocazione a budget finito, operatori di trasferimento $P_1$ e $L^2$, e un benchmark PDE end-to-end su 28 casi held-out;
    \item un'analisi quantitativa che documenta sia il miglioramento di accuratezza osservato con la GNN trajectory-aware sia il costo computazionale e gli interventi del filtro di sicurezza.
\end{enumerate}

Le conclusioni sono limitate alla famiglia di PDE, alla geometria bidimensionale, alla discretizzazione e agli intervalli parametrici effettivamente testati. Non vengono dimostrate generalizzazione a risoluzioni diverse, trasferibilità a PDE differenti, convergenza ottima degli obiettivi di mesh o superiorità in termini di tempo di calcolo. Le prestazioni end-to-end sono attribuite alla pipeline completa GNN--filtro--ALE--FEM e non alla sola rete neurale.

\subsection{Organizzazione del manoscritto}

La Sezione~\ref{sec:problema} definisce il problema di convezione--diffusione, la soluzione manufatta e la discretizzazione. La Sezione~\ref{sec:framework} descrive la rappresentazione grafica, i modelli statico e trajectory-aware e il filtro geometrico. La Sezione~\ref{sec:ale} introduce i trasferimenti di stato e la formulazione ALE con la D-GCL. La Sezione~\ref{sec:protocollo} specifica dataset, protocolli e implementazione. Le Sezioni~\ref{sec:stageabc} e \ref{sec:staged} riportano rispettivamente le campagne preliminari e il confronto PDE end-to-end. Le limitazioni sono discusse nella Sezione~\ref{sec:discussione}; gli approfondimenti matematici e le diagnostiche supplementari sono raccolti nelle appendici.
